\subsection{The arithmetic flag controls the actual lattice minimum}
\label{global-gram-arithmetic-minimum}

Retain one original arithmetic setting.  Let $f=[F:\mathbb Q]$,
$[E:F]=2$, $n=\dim_E V$, $D_F=|\operatorname{disc}F|$ and
$D_E=|\operatorname{disc}E|$.  All vector norms below are the
original Euclidean norms after the coordinate change $P_v^{-1}$;
thus the Hermitian matrix is diagonal with entries $\pm1$.
Use the original $e_a,f_a,z_u$, $0\le a<r$, and put
\[
 E_a=\Bigl(\prod_{v\mid\infty}\|e_{a,v}\|\Bigr)^{1/f},\quad
 F_a=\Bigl(\prod_{v\mid\infty}\|f_{a,v}\|\Bigr)^{1/f},\quad
 Z_u=\Bigl(\prod_{v\mid\infty}\|z_{u,v}\|\Bigr)^{1/f},
\]
\begin{equation}\label{global-gram-minimum-constant}
 \theta_u=\min\bigl(\{1,Z_u^{-1}\}\cup
    \{E_a^{-1},D_F^{-(r-a)/f}F_a^{-1}:0\le a<r\}\bigr)>0,
 \qquad
 \tau_{0,u}=\min\left\{\frac12,
              \frac{\theta_u}{\sqrt n\,D_E^{1/(4f)}}\right\}.
\end{equation}
The sets indexed by $a$ are empty when $r=0$.  These constants
do not depend on the original finite column representative.
They still depend on the actual field and frame.

For an original column $\rho$, write
$a_{u,\rho}=h_{u,\rho}x_u$ and
$\Lambda_{u,\rho}=V(E)\cap a_{u,\rho}L_f$.  Since $x_u$ preserves
$L_f$ and $h_{u,\rho}$ fixes every $e_a,f_a,z_u$, these lattices
are integral for the original Hermitian form and satisfy
\begin{equation}\label{global-gram-minimum-flag}
 e_a,z_u\in\Lambda_{u,\rho},\qquad
 \Lambda_{u,\rho}\cap X=\bigoplus_{a=0}^{r-1}
        \mathfrak d^{a-r}e_a,\qquad X=\langle e_0,\ldots,e_{r-1}\rangle_E.
\end{equation}
Here $\mathfrak d$ is the extension of the different of $F/\mathbb Q$,
so $N_E(\mathfrak d)=D_F^2$.  In particular its fractional powers
are retained; the $f_a$ need not be integral.
These are the literal original lattice and column conventions.
\cite{global-gram-minimum-source}

\begin{theorem}[A column-independent bound for the original minimum]
\label{global-gram-actual-minimum-theorem}
For every original word $u$, column $\rho$, $g\in G_\infty$ and
$k\in K_2$, the source's exact functions satisfy
\begin{equation}\label{global-gram-actual-minimum}
 \operatorname{least}(g,a_{u,\rho}k)
       \ge\frac{\theta_u}{\sqrt n\,D_E^{1/(4f)}}
                      \nu_u(g)^{-1/f},\qquad
 \tau(g,a_{u,\rho}k)\ge\tau_{0,u}\nu_u(g)^{-1/f}.
\end{equation}
Here $\nu_u(g)=\inf_{s\in S_u(F_\infty)}\prod_v\|s_vg_v\|_{\rm op}$
is the original quotient norm.  The function
$\operatorname{least}$ is the infimum of normalized covolume
sizes over \emph{all} nonzero $\mathcal O_E$-submodules of the
original lattice, as in \texttt{leastSize}; it is not replaced
by the shortest-vector length or an infimum over lines.
\end{theorem}

\begin{proof}
First, for every $0\ne y\in\Lambda_{u,\rho}$ and every
$s\in S_u(F_\infty)$ we prove
\begin{equation}\label{global-gram-minimum-product}
                  \prod_v\|s_vy_v\|\ge\theta_u^f.
\end{equation}
The original group acts upper unitriangularly on the ordered
$e$ flag and sends $z_u$ to $z_u+X$.  Its inverse has these
same properties.  There are four exhaustive cases.

If some $\beta_a=h(e_a,y)$ is nonzero, choose the least such
index $a$.  Every $\beta_a$ is in $\mathcal O_E$ by
\eqref{global-gram-minimum-flag}.  At each place, unitarity and
the earlier vanishing pairings give
\[
 h(e_{a,v},s_vy_v)=h(s_v^{-1}e_{a,v},y_v)=\beta_{a,v}.
\]
The diagonal sign matrix has operator norm one.  Consequently
\[
 \prod_v\|s_vy_v\|\ge
   \frac{\prod_v|\beta_{a,v}|}{\prod_v\|e_{a,v}\|}
 =\frac{|N_{E/\mathbb Q}\beta_a|^{1/2}}{E_a^f}
 \ge E_a^{-f}\ge\theta_u^f.
\]

If every $\beta_a$ vanishes, then $y\in X^\perp$.  The number
$b=h(z_u,y)\in\mathcal O_E$ is unchanged by $s$, because
$s^{-1}z_u-z_u\in X$.  If $b\ne0$, the same calculation gives
$\prod_v\|s_vy_v\|\ge Z_u^{-f}\ge\theta_u^f$.

Suppose also $b=0$ but $y\notin X$.  The original signature
and $h(z_u,z_u)<0$ make
$(X^\perp\cap z_u^\perp)/X$ positive definite at every real place.
Thus $c=h(y,y)$ is a nonzero totally positive integer of $F$.
For each place $\|s_vy_v\|^2\ge c_v$, so the product is at
least $N_{F/\mathbb Q}(c)^{1/2}\ge1$.

Finally, if $y\in X$, write $y=\sum_a\alpha_ae_a$, with
$\alpha_a\in\mathfrak d^{a-r}$, and choose the largest index
$a$ with $\alpha_a\ne0$.  Upper unitriangularity preserves this
last coefficient, so $h(f_{a,v},s_vy_v)=\alpha_{a,v}$.
Inclusion of fractional ideals $(\alpha_a)\subseteq\mathfrak d^{a-r}$
gives $|N_{E/\mathbb Q}\alpha_a|\ge D_F^{2(a-r)}$.  Hence
\[
 \prod_v\|s_vy_v\|\ge D_F^{a-r}F_a^{-f}\ge\theta_u^f.
\]
This proves \eqref{global-gram-minimum-product} without choosing
bounded unipotent representatives or changing the metric.

Since $g_v^{-1}y_v=(s_vg_v)^{-1}s_vy_v$, the operator norm
inequality and then the infimum over $s$ yield
\[
 \prod_v\|g_v^{-1}y_v\|\ge\theta_u^f\nu_u(g)^{-1}.
\]
Here $1\le\nu_u(g)<\infty$, because an indefinite unitary matrix
and its inverse have the same operator norm and their product
norm is at least one.  Arithmetic--geometric mean now gives
\begin{equation}\label{global-gram-minimum-shortest}
 \left(\sum_v\|g_v^{-1}y_v\|^2\right)^{1/2}
                 \ge\sqrt f\,\theta_u\nu_u(g)^{-1/f}.
\end{equation}

It remains to pass to the actual higher-rank minimum.  Let
$0\ne R\subseteq\Lambda_{u,\rho}$ be any
$\mathcal O_E$-submodule, of $E$-rank $j$, $1\le j\le n$.
Its original real realization is a lattice in a real space
of dimension $D=2fj$.  It is finitely generated because
$\Lambda_{u,\rho}$ is a lattice over the Noetherian ring
$\mathcal O_E$.  The realization uses real and imaginary
coordinates with the ordinary Euclidean norm, without an
extra $\sqrt2$ factor.  If $\Delta_R$ is its covolume in its
real span, Minkowski's convex body theorem applied to a cube
in that span gives
\[
 \lambda_1(R)\le\sqrt{2fj}\,\Delta_R^{1/(2fj)}.
\]
One can take cubes with side $2(1+\epsilon)\Delta_R^{1/(2fj)}$
and let $\epsilon$ decrease to zero; their vertices have the
displayed Euclidean radius.  The original definition is exactly
\[
 \operatorname{size}(R)
    =\left(2^{fj}D_E^{-j/2}\Delta_R\right)^{1/(2fj)}
    =\sqrt2\,D_E^{-1/(4f)}\Delta_R^{1/(2fj)}.
\]
Together with \eqref{global-gram-minimum-shortest}, this yields
\[
 \operatorname{size}(R)\ge
     \frac{\theta_u}{\sqrt j\,D_E^{1/(4f)}}\nu_u(g)^{-1/f}
 \ge\frac{\theta_u}{\sqrt n\,D_E^{1/(4f)}}\nu_u(g)^{-1/f}.
\]
The source's $\operatorname{lamX}(g,a_{u,\rho}k)$ is precisely
$\Lambda_{u,\rho}$, since $k\in K_2$ preserves $L_f$; its
\texttt{frm} inverse is precisely $g_v^{-1}P_v^{-1}$, or its
complex conjugate at the other representative of a complex
place.  Thus these are the actual covolumes in
\texttt{leastSize}.  Taking its infimum proves the first
assertion.  The second follows from $\nu_u\ge1$ and
$\tau=\min\{1/2,\operatorname{least}\}$.
\cite{global-gram-minimum-source}
\end{proof}

\paragraph{Exact effect on the projector's far-return bound.}
Retain the original projector kernel $\phi$, its finite idempotent,
formal degrees and positive normalizer
$\mathcal M_\phi=(\prod_vd_v)e(1)>0$.
For a cutoff $\varepsilon$, write
\[
 \mathcal T_{h,\varepsilon}(x)=
    \mathcal M_\phi^{-1}\sum_{\gamma\in G(F)\text{ far at }x}
                                  |\phi(x^{-1}\gamma x)|.
\]
Assume the literal source \texttt{Deep.FarBound}, with its
discrete-series, signature and gap hypotheses and the conditions
\[
 \begin{gathered}
 q=\min(p,q)\ge1,\quad0<\varepsilon\le(2q)^{-1/2},\quad
  2n^2\varepsilon\le\log2,\\
 64n^2\le h\varepsilon,\quad
 1\le h\varepsilon^2/16,\quad120nq^{3/2}\le h.
 \end{gathered}
\]
That source estimate has authenticated interfaces; its entire
dependency chain is not reproved here.  Substitution into it gives
the actual pointwise bound, uniformly over the original columns
and $k\in K_2$,
\begin{equation}\label{global-gram-minimum-far}
 \mathcal T_{h,\varepsilon}(g,a_{u,\rho}k)
 \le C_u e^{-h\varepsilon^2/16}\nu_u(g)^{4n^2},\qquad
 C_u=2f3^f\left(\frac{3\sqrt n}{\tau_{0,u}^2}\right)^{2n^2f}.
\end{equation}
The exponent is exactly $4n^2$, since the original factor is
$\tau^{-4n^2f}$ and \eqref{global-gram-actual-minimum} has
$\nu^{-1/f}$; no degree has been suppressed.

In the original seed-decay gap regime, set
\[
 \mathcal W_u=\int_{S_u\backslash G_\infty}
                \nu_u(g)^{4n^2}\|W_u(g)\|\,dq_u(g)<\infty.
\]
Finiteness follows from Theorem~\ref{actual-seed-decay-theorem},
whose exponent is $N_X\ge4n^2$, and $\nu_u\ge1$.
Using the original representative measure $\nu_Q(u)$ gives
\begin{equation}\label{global-gram-minimum-weighted-far}
 \int_{G_\infty}\mathcal T_{h,\varepsilon}(g,a_{u,\rho}k)
                      \|W_u(g)\|\,d\nu_Q(u)(g)
          \le C_u e^{-h\varepsilon^2/16}\mathcal W_u<\infty.
\end{equation}
The integrand's far-return factor need not descend to the
$S_u$ quotient; the displayed measure is therefore kept
explicit.  Its majorant does descend, and the established
Weil comparison gives the stated right side.  At the already
admissible choice $h\ge16n^{19}$ and
$\varepsilon=(8n^2)^{-1}$ the exponential is at most
$e^{-n^{15}/64}$.\cite{global-gram-minimum-source}

This supplies an arithmetic minimum estimate and a finite weighted
far-return majorant on the actual period cells.  A bound for
$C_u\mathcal W_u$ across the moving arithmetic family, control
of the near returns, and identification with the original
restricted arithmetic chain and its positive trace remain required
for the numerical budget.  None is a consequence of finiteness
in \eqref{global-gram-minimum-weighted-far}.
